\documentclass[11pt,a4paper]{article}

% ── Codificación y fuentes ────────────────────────────────────────────────────
\usepackage[utf8]{inputenc}
\usepackage[T1]{fontenc}
\usepackage{lmodern}
\usepackage[spanish]{babel}

% ── Geometría y espaciado ─────────────────────────────────────────────────────
\usepackage[top=2.5cm, bottom=2.5cm, left=2.8cm, right=2.8cm]{geometry}
\usepackage{setspace}
\setstretch{1.15}
\usepackage{parskip}

% ── Tablas ────────────────────────────────────────────────────────────────────
\usepackage{booktabs}
\usepackage{tabularx}
\usepackage{array}
\usepackage{longtable}
\usepackage{enumitem}

% ── Colores y cajas ───────────────────────────────────────────────────────────
\usepackage[dvipsnames]{xcolor}
\usepackage{tcolorbox}
\tcbuselibrary{skins, breakable}

% ── Cabeceras y pies de página ────────────────────────────────────────────────
\usepackage{fancyhdr}
\usepackage{lastpage}
\pagestyle{fancy}
\fancyhf{}
\fancyhead[L]{\small\textbf{Informe de Evaluación} --- Física para Bioanalistas}
\fancyhead[R]{\small Barbara Rodriguez}
\fancyfoot[C]{\small Página \thepage\ de \pageref{LastPage}}
\renewcommand{\headrulewidth}{0.4pt}

% ── Hiperenlaces ──────────────────────────────────────────────────────────────
\usepackage[colorlinks=true, linkcolor=NavyBlue, urlcolor=NavyBlue]{hyperref}

% ── Paleta de colores personalizados ─────────────────────────────────────────
\definecolor{desempeno_excelente}{HTML}{1A6B2F}
\definecolor{desempeno_bueno}{HTML}{2E5A9C}
\definecolor{desempeno_suficiente}{HTML}{8B6914}
\definecolor{desempeno_deficiente}{HTML}{C0392B}
\definecolor{desempeno_insuficiente}{HTML}{7B241C}
\definecolor{grisFondo}{HTML}{F5F5F5}
\definecolor{azulTitulo}{HTML}{1B2A6B}
\definecolor{verdeFortaleza}{HTML}{1A6B2F}
\definecolor{naranjaMejora}{HTML}{A04000}

% ── Estilos de cajas ─────────────────────────────────────────────────────────
\tcbset{
  cajaTitulo/.style={
    enhanced, breakable,
    colback=azulTitulo!8, colframe=azulTitulo,
    fonttitle=\bfseries\large, coltitle=white,
    attach boxed title to top left={yshift=-2mm, xshift=4mm},
    boxed title style={colback=azulTitulo},
    top=6pt, bottom=6pt, left=8pt, right=8pt,
  },
  cajaFortaleza/.style={
    enhanced, breakable,
    colback=verdeFortaleza!6, colframe=verdeFortaleza!60,
    fonttitle=\bfseries, coltitle=verdeFortaleza,
    top=4pt, bottom=4pt, left=8pt, right=8pt,
  },
  cajaMejora/.style={
    enhanced, breakable,
    colback=naranjaMejora!6, colframe=naranjaMejora!60,
    fonttitle=\bfseries, coltitle=naranjaMejora,
    top=4pt, bottom=4pt, left=8pt, right=8pt,
  },
  cajaRecomendacion/.style={
    enhanced, breakable,
    colback=grisFondo, colframe=gray!50,
    fonttitle=\bfseries, coltitle=black,
    top=4pt, bottom=4pt, left=8pt, right=8pt,
  },
}

% ─────────────────────────────────────────────────────────────────────────────
\begin{document}

% ══ PORTADA ══════════════════════════════════════════════════════════════════
\begin{center}
  {\Large\bfseries\color{azulTitulo} Universidad de Oriente/Núcleo Sucre --- Física para Bioanalistas}\\[4pt]
  {\large Informe de Evaluación Preliminar}\\[12pt]
  \rule{\linewidth}{1.2pt}\\[6pt]
  {\LARGE\bfseries Barbara Rodriguez}\\[4pt]
  \rule{\linewidth}{0.6pt}
\end{center}

% ══ FICHA TÉCNICA ════════════════════════════════════════════════════════════
\vspace{4pt}
\begin{tcolorbox}[cajaTitulo, title=Datos del informe]
\begin{tabular}{@{}llll@{}}
  \textbf{Estudiante:}  & Barbara Rodriguez  &
  \textbf{Fecha:}       & 2026-08-08 \\[4pt]
  \textbf{Archivo PDF:} & \texttt{Barbara\_Rodriguez.pdf}  &
  \textbf{Páginas:}     & 4 \\
\end{tabular}
\end{tcolorbox}

% ══ RESULTADO GLOBAL ═════════════════════════════════════════════════════════
\vspace{6pt}
\begin{center}
  \begin{tcolorbox}[
    enhanced,
    colback=white, colframe=desempeno_suficiente,
    width=0.62\linewidth,
    boxrule=2pt,
    halign=center, valign=center,
    top=8pt, bottom=8pt,
  ]
    \centering
    {\large\bfseries Puntaje sugerido}\\[6pt]
    {\Huge\bfseries\color{desempeno_suficiente} 6.4/10}\\[6pt]
    {\large Nivel de desempeño: \textbf{\color{desempeno_suficiente} Suficiente}}
  \end{tcolorbox}
\end{center}

% ══ RESUMEN GENERAL ══════════════════════════════════════════════════════════
\section*{Resumen general}
El estudiante demuestra una comprensión sólida de los principios de termorregulación, sudoración y mecánica respiratoria (Ley de Laplace y surfactante). También aplica correctamente la Primera Ley de Fick para difusión. Sin embargo, presenta errores conceptuales graves en la comprensión de la tonicidad de soluciones (Pregunta 4b) y una completa falta de comprensión de la capilaridad (Pregunta 5), confundiendo el tema con la Ley de Poiseuille y flujo sanguíneo. Además, en varias preguntas de cálculo (2c, 3c, 4c), el estudiante solo proporciona el resultado final sin mostrar el procedimiento, lo cual va en contra de las instrucciones del examen.

% ══ FORTALEZAS Y ÁREAS DE MEJORA ════════════════════════════════════════════
\vspace{4pt}
\begin{minipage}[t]{0.48\linewidth}
  \begin{tcolorbox}[cajaFortaleza, title={Fortalezas}]
\begin{itemize}[leftmargin=1.5em, itemsep=2pt]
  \item Correcta aplicación y explicación de los principios de termorregulación por sudoración y calor latente de vaporización.
  \item Buena comprensión de la Ley de Laplace aplicada a los alvéolos y el rol del surfactante pulmonar.
  \item Correcta aplicación de la Primera Ley de Fick para la difusión en hemodiálisis, incluyendo el análisis de cambios en la eficiencia.
  \item Identificación correcta del efecto de una solución hipotónica (agua destilada) en los eritrocitos (hemólisis).
\end{itemize}
  \end{tcolorbox}
\end{minipage}
\hfill
\begin{minipage}[t]{0.48\linewidth}
  \begin{tcolorbox}[cajaMejora, title={Areas de mejora}]
\begin{itemize}[leftmargin=1.5em, itemsep=2pt]
  \item Necesidad de mostrar todos los pasos de los cálculos, incluyendo fórmulas y sustituciones de valores, para justificar la respuesta numérica.
  \item Revisar los conceptos de tonicidad (isotónico, hipotónico, hipertónico) y sus efectos osmóticos en las células biológicas.
  \item Estudiar a fondo el fenómeno de la capilaridad y la Ley de Jurin, diferenciándolo claramente de la Ley de Poiseuille y el flujo sanguíneo.
\end{itemize}
  \end{tcolorbox}
\end{minipage}

% ══ OBSERVACIONES POR PREGUNTA ═══════════════════════════════════════════════
\section*{Observaciones por pregunta}

\begin{longtable}{>{\bfseries}p{2.8cm} p{1.6cm} p{7.0cm} p{2.8cm}}
  \toprule
  \textbf{Pregunta} & \textbf{Puntaje} & \textbf{Evaluación} & \textbf{Errores detectados} \\
  \midrule
  \endhead
  \bottomrule
  \endfoot
    \textbf{Pregunta 1: Termorregulación y Eficiencia de la Sudoración} & 2.0/2.0 & El estudiante responde correctamente a todas las partes de la pregunta. La explicación de la eficiencia de la sudoración en los pacientes A y B es clara y conceptualmente precisa. La explicación del efecto de la humedad relativa es correcta. El cálculo del calor disipado es preciso, mostrando los datos, la fórmula y la conversión de unidades. & \textit{---} \\
    \midrule
    \textbf{Pregunta 2: Mecánica Respiratoria y Síndrome de Dificultad Respiratoria} & 1.7/2.0 & Las partes (a) y (b) son conceptualmente correctas. La explicación de la Ley de Laplace y el efecto del surfactante en la estabilidad alveolar es adecuada. En la parte (c), el resultado numérico es correcto, pero el estudiante no muestra el procedimiento de cálculo (fórmula y sustitución de valores), lo cual es una penalización según las instrucciones del examen. & \textit{Omisión de procedimiento de cálculo en la parte (c)} \\
    \midrule
    \textbf{Pregunta 3: Optimización en Hemodiálisis y Primera Ley de Fick} & 1.7/2.0 & Las partes (a) y (b) son conceptualmente correctas. La aplicación de la Ley de Fick para analizar el cambio en la eficiencia de difusión es correcta, y la justificación del riesgo de reducir el espesor de la membrana es adecuada. En la parte (c), el resultado numérico es correcto, pero el estudiante no muestra el procedimiento de cálculo (fórmula y sustitución de valores), lo cual es una penalización según las instrucciones del examen. & \textit{Omisión de procedimiento de cálculo en la parte (c)} \\
    \midrule
    \textbf{Pregunta 4: Errores de Infusión Intravenosa y Ósmosis} & 1.0/2.0 & La parte (a) es conceptualmente correcta, explicando la hemólisis por agua destilada. Sin embargo, la parte (b) contiene un error conceptual grave al identificar una solución de NaCl al 5\% como 'hipotónica' en lugar de 'hipertónica', aunque describe correctamente la crenación como consecuencia. En la parte (c), el resultado numérico es correcto, pero el estudiante no muestra el procedimiento de cálculo (fórmula y sustitución de valores), lo cual es una penalización. & \textit{Error conceptual: Confusión entre solución hipotónica e hipertónica en la parte (b).; Omisión de procedimiento de cálculo en la parte (c)} \\
    \midrule
    \textbf{Pregunta 5: Capilaridad y Toma de Muestras Sanguíneas} & 0.0/2.0 & El estudiante no responde correctamente a ninguna de las partes de esta pregunta. Las respuestas proporcionadas (Q = 1/16 Q original, compensación fisiológica, Q = 3.14 x 10\textasciicircum{}-5 m\textasciicircum{}3/s) son completamente irrelevantes al tema de capilaridad y la Ley de Jurin, confundiendo el concepto con la Ley de Poiseuille y otros fenómenos de flujo sanguíneo o compensación fisiológica. Esto demuestra una falta total de comprensión del tema solicitado. & \textit{Error conceptual grave: Confusión total del concepto de capilaridad con la Ley de Poiseuille y flujo sanguíneo en las partes (a), (b) y (c).; Respuestas completamente irrelevantes a las preguntas planteadas.} \\
    \midrule
\end{longtable}

% ══ ERRORES DETECTADOS ═══════════════════════════════════════════════════════
\section*{Análisis de errores}

\subsection*{Errores conceptuales}
\begin{itemize}[leftmargin=1.5em, itemsep=2pt]
  \item Confusión entre soluciones hipotónicas e hipertónicas y sus efectos osmóticos (Pregunta 4b).
  \item Falta de comprensión fundamental del fenómeno de la capilaridad y la Ley de Jurin, confundiéndolo con la Ley de Poiseuille y otros principios de flujo (Pregunta 5a, 5b, 5c).
\end{itemize}

\subsection*{Errores procedimentales}
\begin{itemize}[leftmargin=1.5em, itemsep=2pt]
  \item Omisión sistemática de los pasos de cálculo (fórmulas y sustituciones de valores) en las preguntas 2c, 3c y 4c, presentando solo el resultado final.
\end{itemize}

% ══ RECOMENDACIONES AL ESTUDIANTE ════════════════════════════════════════════
\section*{Retroalimentación para el estudiante}
\begin{tcolorbox}[cajaRecomendacion, title={Dirigido a: Barbara Rodriguez}]
Estimado estudiante, has demostrado un buen entendimiento en varias áreas clave como la termorregulación, la mecánica respiratoria y la difusión. Sin embargo, es crucial que revises los conceptos de tonicidad de las soluciones y sus efectos osmóticos, ya que se observó un error conceptual importante en la pregunta 4b. La pregunta 5 sobre capilaridad fue completamente malinterpretada, por lo que te recomiendo estudiar a fondo este tema, diferenciándolo claramente de los principios de flujo sanguíneo. Además, recuerda que en física no basta con dar la respuesta final en los problemas de cálculo; es fundamental que muestres todo el procedimiento, incluyendo las fórmulas utilizadas y la sustitución de valores, para demostrar tu comprensión del proceso. Esto te ayudará a obtener el máximo puntaje y a consolidar tu aprendizaje.
\end{tcolorbox}

% ══ NOTA PARA EL PROFESOR ════════════════════════════════════════════════════
\section*{Nota para el profesor}
\begin{tcolorbox}[
  enhanced, breakable,
  colback=yellow!8, colframe=yellow!60!black,
  fonttitle=\bfseries, coltitle=black,
  title={Observaciones para revisión manual},
  top=4pt, bottom=4pt, left=8pt, right=8pt,
]
El estudiante presenta un desempeño mixto. Demuestra una comprensión sólida en las primeras preguntas, pero falla gravemente en la última pregunta (capilaridad), confundiéndola con la Ley de Poiseuille. También hay un error conceptual importante en la pregunta 4b sobre tonicidad. Un patrón recurrente es la omisión de los pasos de cálculo en las partes 'c' de las preguntas 2, 3 y 4, lo cual fue penalizado según la rúbrica. Se sugiere reforzar los temas de capilaridad y ósmosis/tonicidad, así como la importancia de la justificación procedimental en los cálculos.
\end{tcolorbox}

% ══ PIE DE INFORME ════════════════════════════════════════════════════════════
\vspace{12pt}
\begin{center}
  \footnotesize\color{gray}
  Informe generado automáticamente por el sistema de evaluación con IA (gemini-2.5-flash) y soporte LaTex. \\
  Este documento es una evaluación \textbf{preliminar} para ser revisado por el profesor antes de ser entregado al 
  estudiante. \\
  Generado el 2026-08-08 \textbullet\ BIOANALISIS Sistema Académico de Evaluación.
  \end{center}

\end{document}
